% Test approach.

% {method}{yes / no}{remark}
\VerificationMethod{Requirements-based test}{yes}{}
\VerificationMethod{Interface test}{yes}{}
\VerificationMethod{Fault injection test}{yes}{NvM failures}
\VerificationMethod{Resource usage evaluation}{no}{done at integration level}
\VerificationMethod{Back-to-back test (model vs.\ code)}{no}{no model}

% one line per technique
\DerivationTechnique{analysis of requirements and detailed design}
\DerivationTechnique{equivalence class partitioning of inputs and states}
\DerivationTechnique{boundary value analysis (min, max, max+1)}
\DerivationTechnique{error guessing}

% {metric}{target}{justification if not met} -- plain numbers get a % sign
\Coverage{Statement}{100}{review of uncovered lines}
\Coverage{Branch}{100}{review of uncovered branches}
\Coverage{MC/DC}{n.a.}{not required for ASIL B}

\EntryCriterion{Unit source is baselined (\TheBaseline) and compiles without warnings.}
\EntryCriterion{Requirements \DocRef{srs} and design \DocRef{design} are reviewed.}
\ExitCriterion{All test cases executed; all passed or deviations accepted.}
\ExitCriterion{Coverage targets reached or gaps justified.}

\PassFailCriteria{
  A test case passes if every step yields the expected result and no
  unexpected side effect (assertion, sanitizer report, timeout) occurs.
}

% Free text example: shared figure code from the root figures/ folder (used by all
% units), then an own subsection with a self-made table and interval notation.
\begin{FreeText}
\subsection{Test Process}
\begin{center}
\InputFigure{test_process}
\end{center}

\subsection{Equivalence Classes of the Setpoint}
\begin{center}
\begin{tabular}{lll}
  \toprule
  \textbf{Class} & \textbf{Range} & \textbf{Representatives (boundaries)} \\ \midrule
  invalid low  & $sp < 0$            & not representable (\texttt{uint8}) \\
  valid        & $0 \le sp \le 100$  & 0, 1, 50, 99, 100 \\
  invalid high & $100 < sp \le 255$  & 101, 255 \\
  \bottomrule
\end{tabular}
\end{center}
Statement coverage is measured as
\[
  C_{\mathrm{stmt}} = \frac{\text{executed statements}}{\text{all statements}} \cdot 100\,\% .
\]
\end{FreeText}
